% Research notes for incorporation into the main article.
% This is a LaTeX fragment: the parent manuscript supplies amsmath/amsthm.
% Bibliographic records at the end may be merged into its bibliography.

\section{Repetition kernels from aligned component partitions}
\label{sec:aligned-routing}

The four-copy theorem for connected tiles extends to a class of disconnected
tiles substantially larger than matchings of individual input and output ports.
The correct hypothesis concerns the partitions of the interface indices induced
by connected components.  In particular, componentwise transport avoids taking
the least common multiple of permutation cycle lengths.

\subsection{Conventions and the connected primitive}

All graphs are finite undirected multigraphs; loops and parallel edge occurrences
are retained.  Every edge has unit cut cost.  A tile has ordered disjoint
interfaces
\[
 L=(\ell_1,\ldots,\ell_w),\qquad R=(r_1,\ldots,r_w),
\]
where every listed vertex is distinct and $w\geq1$.  Serial composition identifies
the output of one occurrence with the input of the next, coordinate by coordinate.
It identifies vertices, never edges.  Private vertices in distinct occurrences
remain distinct unless a prescribed serial seam identifies boundary vertices.

For a coloring $c:V(G)\to\{0,1\}$ write
\[
 \begin{aligned}
 q_G(c)&=|\{uv\in E(G):c(u)\ne c(v)\}|,\\
 F_G(c)&=\{a:\text{an edge has both endpoints colored }a\}.
 \end{aligned}
\]
For a threshold $k\geq0$, the marked signature $\mathsf S_k(G)$ is the minimum
cut cost for each prescribed coloring of its exposed vertices and each
\emph{exact} mark set $F\subseteq\{0,1\}$, with every value exceeding $k$ and
infeasibility represented by the common value $k+1$.

The incoming four-copy article proves the following primitive
\cite{incoming-four-copy}.  It is recalled to make the dependency of the new
result explicit.

\begin{lemma}[Connected repetition primitive]
\label{lem:route-connected}
For a connected tile $W$ with equal nonempty input and output widths, put
\[
 a_k=\begin{cases}4,&k=3,\\ k+2,&k\ne3.\end{cases}
\]
Then, for every positive integer $q$,
\[
 \mathsf S_k(W^{\circ q})
 =\mathsf S_k(W^{\circ\min(q,a_k)}).
\]
Moreover, a coloring of $W^{\circ a_k}$ with cut cost at most $k$ contains a
zero-cost occurrence of $W$.  Such an occurrence is monochromatic and may be
duplicated arbitrarily without changing boundary colors, cut cost, or exact
edge marks.
\end{lemma}

The duplication assertion follows because $a_k>k$, every positive-cost
occurrence contributes at least one, and connectedness makes a zero-cost
coloring monochromatic.  The distinction between equality of \emph{minimum}
signatures and preservation of every individual coloring remains relevant:
compression in the connected four-copy proof can lower an individual coloring's
cost; lifting from the kernel by duplication preserves its cost exactly.

\subsection{Aligned partitions}

\begin{definition}[Aligned component partitions]
\label{def:route-aligned}
Let $C_1,\ldots,C_d$ be all connected components of a tile $T$.  Define
\[
 A_j=\{i:\ell_i\in C_j\},\qquad B_j=\{i:r_i\in C_j\}.
\]
The tile has aligned component partitions if every $A_j$ and $B_j$ is nonempty
and the two unordered collections of subsets of $[w]$ coincide:
\[
 \{A_1,\ldots,A_d\}=\{B_1,\ldots,B_d\}.
\]
There is then a unique permutation $\sigma$ of $[d]$ satisfying
$B_j=A_{\sigma(j)}$ for every $j$.
\end{definition}

The sets in either collection are a partition of $[w]$.  The condition is
testable directly from the component graph and the port lists.  It includes a
connected tile ($d=1$), and the permutation-zero class where every component
has exactly one input and one output ($d=w$).  It also permits variable component
widths.  For example, one component can transport a one-element block to a
two-element block, provided another component transports the latter block back.
Thus $|A_j|=|B_j|$ is \emph{not} required for individual components.

Within component $C_j$, order its input ports by their original indices in
$A_j$ and its output ports by their original indices in $A_{\sigma(j)}$.
Its individual input and output widths may differ.  Adjacent components in a
$\sigma$-orbit have matching intermediate interfaces because they use exactly
the same index set.  A complete orbit word returns to its initial index set
and therefore is an ordinary equal-width connected tile.

\begin{lemma}[Decomposition into component chains]
\label{lem:route-chains}
Suppose $T$ has aligned component partitions.  For every $N\geq1$, its open
power is the disjoint union, over $j\in[d]$, of the connected component chains
\[
 P_{j,N}=C_j\circ C_{\sigma(j)}\circ\cdots\circ C_{\sigma^{N-1}(j)}.
\]
The input ports of this chain are $L|_{A_j}$ and its output ports are
$R|_{A_{\sigma^N(j)}}$, with their original coordinate labels.
\end{lemma}

\begin{proof}
At a serial seam, all outputs of $C_j$ have indices in $B_j=A_{\sigma(j)}$,
so they are identified with inputs of exactly one component in the next layer,
namely $C_{\sigma(j)}$.  Every one of those input ports is used.  Conversely,
each component in the next layer has exactly one predecessor component because
$\sigma$ is a permutation.  Consequently distinct chains do not merge or split.
Every component occurrence is on exactly one chain, the one whose initial
index is $\sigma^{-s}(j)$ if it is $C_j$ in layer $s$.  Each chain is connected
because its components are connected and consecutive components are glued
along a nonempty interface.  The assertion about external ports follows by
tracking these index sets through the $N$ layers.
\end{proof}

\subsection{The componentwise kernel theorem}

Let $\mathcal O$ run over the cycles of $\sigma$, and put
\[
 h_{\mathcal O}=|\mathcal O|,\qquad
 Q_{\mathcal O}=\sum_{j\in\mathcal O}
       \bigl(|V(C_j)|+|E(C_j)|\bigr),\qquad
 Q=\sum_{\mathcal O}Q_{\mathcal O}.
\]
For a fixed $j\in\mathcal O$, define the full-cycle word and prefix word by
\[
 W_j=C_j\circ\cdots\circ C_{\sigma^{h_{\mathcal O}-1}(j)},
 \qquad
 V_{j,r}=C_j\circ\cdots\circ C_{\sigma^{r-1}(j)}.
\]
The word $V_{j,0}$ is empty: it is an identity interface and contributes no
new edge or private vertex.  This convention is used in formulas; an
implementation need not create a zero-exponent block.

\begin{theorem}[Kernel without a global permutation period]
\label{thm:route-kernel}
Let $T$ have aligned component partitions and let $N\geq1$ be given in binary.
For each cycle $\mathcal O$, write
\[
 N=q_{\mathcal O}h_{\mathcal O}+r_{\mathcal O},
 \qquad 0\leq r_{\mathcal O}<h_{\mathcal O}.
\]
Replace the chain starting at $j\in\mathcal O$ by
\[
 K_{j,N}
 =W_j^{\circ\min(q_{\mathcal O},a_k)}
      \circ V_{j,r_{\mathcal O}},
\]
omitting an empty factor.  Give this chain the original exposed input labels
$A_j$ and output labels $A_{\sigma^N(j)}$.  Let $K(T,N)$ be the disjoint union
of these $d$ chains.  Then
\[
 \mathsf S_k(T^{\circ N})=\mathsf S_k(K(T,N)).
\]
Its explicit size, counting vertices and edge occurrences before any external
gluing, is at most
\[
 (a_k+1)\sum_{\mathcal O}h_{\mathcal O}Q_{\mathcal O}
 \leq (a_k+1)dQ.
\]
In particular the bound is $5\sum_{\mathcal O}h_{\mathcal O}Q_{\mathcal O}
\leq5dQ$ at threshold three, and is bounded by
$(k+3)\sum_{\mathcal O}h_{\mathcal O}Q_{\mathcal O}$ for every $k$.
For fixed $k$ the kernel is constructible in polynomial bit time in the explicit
input size and $\log(N+1)$.  No global order $\operatorname{ord}(\sigma)$ is
constructed as a word length or used for expansion.
\end{theorem}

\begin{proof}
By Lemma~\ref{lem:route-chains},
\[
 P_{j,N}=W_j^{\circ q_{\mathcal O}}\circ V_{j,r_{\mathcal O}}.
\]
The full-cycle word $W_j$ is connected and has input and output index set
$A_j$, so Lemma~\ref{lem:route-connected} applies to its repeated factor when
$q_{\mathcal O}>0$.  If $q_{\mathcal O}=0$, then $N<h_{\mathcal O}$ and the
entire chain is the explicitly retained prefix of $N$ components.  There is
therefore no use of the connected primitive at exponent zero.

Composition with the unchanged prefix preserves the marked signature:
compatible boundary colors agree on the seam, cut costs add because edges
are disjoint, and exact marks combine by union.  The same argument for
disjoint union combines the $d$ chain replacements.  Saturation at $k+1$ is
legitimate because all cut costs are nonnegative.  This proves signature
equality with all original exposed colors prescribed.

There is a potentially important boundary-label issue.  The retained number
of component occurrences is
$\min(q_{\mathcal O},a_k)h_{\mathcal O}+r_{\mathcal O}$, which is congruent
to $N$ modulo $h_{\mathcal O}$.  Thus the terminal component index is the
original one, $\sigma^N(j)$.  More importantly, its actual port labels are
retained coordinatewise.  Different cycles may have different retained
lengths; the output of the construction is not asserted to be a common
power $T^{\circ M}$ for one integer $M$.

Each full-cycle word contains one occurrence of each component in
$\mathcal O$, so its size is at most $Q_{\mathcal O}$.  The prefix has size
at most $Q_{\mathcal O}$.  Each of the $h_{\mathcal O}$ chains associated
with the cycle therefore contributes at most $(a_k+1)Q_{\mathcal O}$.
Summing gives the stated bound; vertex identifications only reduce it.

Connected components, index partitions, and permutation cycles are computed
from the explicit tile.  Division of a binary $N$ by an integer at most $d$,
followed by comparison of its quotient with the fixed number $a_k$, is
polynomial in the encoded input size.  The construction writes only the
bounded number of component occurrences just counted.  If graph identifiers
are encoded in binary, the bit size adds the usual logarithmic factor for
identifiers; it remains polynomial.  This also proves the complexity claim.
\end{proof}

\begin{corollary}[Singleton routes]
If every component has one input and one output, write $\pi(i)$ for the
output index of the component containing $\ell_i$.  Then $T^{\circ N}$ has
$w$ connected two-terminal chains, the chain starting at $\ell_i$ ending at
$r_{\pi^N(i)}$.  Theorem~\ref{thm:route-kernel} gives a threshold-three
kernel of size at most $5wQ$ using the individual cycle lengths of $\pi$.
\end{corollary}

\subsection{Contextual replacement and mixed flat descriptions}

\begin{theorem}[Contextual equivalence]
\label{thm:route-context}
Suppose an arbitrary finite graph context contains edge-disjoint tile-power
macros.  Each macro has either a connected tile or aligned component
partitions.  The context may identify any of the exposed macro vertices
with each other or with context vertices, including ports belonging to the
same macro.  It has no access to private macro vertices.

Replacing every macro by its kernel preserves the marked signature at
threshold $k$ on any exposed context vertices.  In particular it preserves
the existence of a cut of size at most $k$ whose two shores each contain an
edge.  Connected macros contribute $O(Q_i)$ at fixed $k$; an aligned macro
contributes
$O(\sum_{\mathcal O}h_{\mathcal O}Q_{\mathcal O})=O(d_iQ_i)$.
Thus any finite flat mixture of these macros has a polynomial-size explicit
query kernel, independently of the numeric multiplicities.
\end{theorem}

\begin{proof}
Fix a coloring of the context vertices, including every equivalence class
of identified exposed vertices.  Its restriction prescribes compatible
colors on the exposed vertices of every macro.  For each choice of exact
local mark sets, the optimal global cost is the sum of the context cost
and the optimal macro costs; the global mark set is the union of their
exact mark sets.  Theorem~\ref{thm:route-kernel} and the connected primitive
preserve each summand at threshold $k$.  Taking minima over local mark
sets and all allowed context colorings proves the assertion.

The argument permits arbitrary identifications among exposed vertices:
they merely restrict the allowed colorings.  It does not permit an
additional edge or gluing at a private vertex, because the signature did
not record the color of that vertex.  All sizes add over the explicitly
listed macros and context, and quotienting can only reduce vertex count.
\end{proof}

This theorem concerns a \emph{flat} input: a capped tile is explicit.
Unrestricted nesting of macro definitions is not covered by the size
bound without an additional argument controlling intermediate expansion.
The bound also does not by itself validate a spherical embedding, preserve
the graph's entire topology, or certify that an ambient 3-manifold is a
ball.  Its topological use is the threshold-three terminal query on an
actual spherical cubic pattern in a separately certified ball.

\subsection{Exact source-witness lifting}

\begin{theorem}[Cycle-word run certificates]
\label{thm:route-lift}
A threshold-three marked cut of $K(T,N)$ lifts to a coloring of the original
$T^{\circ N}$ with the same cost, exposed colors, and exact marks.
For each component chain, its lifted description uses at most four retained
full-cycle coloring blocks and one prefix coloring block.  One of the
full-cycle blocks may carry a binary repetition count.  The description and
its verification are polynomial in the encoded source and kernel sizes.
At most three original cut-edge addresses and two original monochromatic
edge addresses suffice to list its cut and its two edge witnesses.
\end{theorem}

\begin{proof}
If $q_{\mathcal O}\leq4$, the chain was not shortened, and its explicit
kernel coloring is already a source coloring.  If $q_{\mathcal O}>4$,
its kernel contains four copies of the connected word $W_j$.  Their total
cost is at most the global cost, at most three.  One of these four copies
therefore has zero cost and is monochromatic.  Replace it by
$q_{\mathcal O}-3$ copies with that same monochromatic coloring.  Its
input and output colors agree, so the seams remain valid.  Its cost remains
zero, and its monochromatic edge mark is unchanged.  All other full-cycle
blocks and the prefix are left unchanged.

Do this independently on each component chain.  The chains have no internal
shared vertices, and their exposed colors are unchanged, so every external
identification remains valid.  Exact global marks and cost are preserved.
This produces the asserted run representation.

An edge in component $C_{\sigma^u(j)}$ of full-cycle occurrence $v$ has
original tile-copy index $vh_{\mathcal O}+u$, using zero-based indices.
An edge in the prefix has original tile-copy index
$q_{\mathcal O}h_{\mathcal O}+u$.  Keep its original local edge-occurrence
identifier.  These formulas locate edges in the source without enumerating
omitted copies.  A positive-cost run can contain at most three cut edges
in total, while a repeated monochromatic run contributes none.  Consequently
all cut addresses can be listed directly, and a first monochromatic edge
for each mark gives the two mark witnesses.  The copy indices have
$O(\log(N+1))$ bits.  Checking local colors, seams, edge costs, mark unions,
and the binary run counts requires only polynomial work.
\end{proof}

For a general fixed threshold $k$, use $a_k$ full-cycle blocks and the
same proof, since $a_k>k$.  The two edge-mark witnesses remain sufficient;
the cut list has at most $k$ addresses.

\subsection{Hypotheses that must survive implementation}

\paragraph{Original output transport.}
Consider the width-two matching with edges $\ell_1r_2$ and $\ell_2r_1$.
At exponent five, input colors $(0,1)$ and output colors $(1,0)$ admit
a zero-cost coloring with both marks.  Blindly replacing the exponent
by four changes the output transport to the identity; those boundary
colors then require positive cost.  Cyclewise reduction preserves the
original labeled output, whereas a single global cap need not.

\paragraph{A partition mismatch really can merge chains.}
Let a width-three tile have two components, one with edges
$\ell_1r_1,\ell_2r_1$, the other with edges
$\ell_3r_2,\ell_3r_3$.  Its input partition is
$\{\{1,2\},\{3\}\}$ while its output partition is
$\{\{1\},\{2,3\}\}$.  The first tile occurrence has two components,
but the square of the tile is connected: the second occurrence's component
containing $\ell_1,\ell_2$ joins the two previous components, and the other
second-layer component attaches to the same union.  Thus there is no fixed
two-chain decomposition.  This example disproves that decomposition outside
the alignment condition; it does not prove that every possible kernel fails.
Section~\ref{sec:through-reblocking} supplies a verified reblocking strategy
that applies to this example and, more generally, to every tile whose
components all meet both interfaces.

\paragraph{Naive independent strands lose constraints.}
A connected four-spoke star with boundary leaves
$\ell_1,\ell_2,r_1,r_2$ cannot be replaced by two independent paths
merely by pairing $\ell_i$ with $r_i$.  Prescribing
$c(\ell_1)=c(r_1)=0$ and $c(\ell_2)=c(r_2)=1$ gives zero cost and both
marks on the two independent paths, but is not a zero-cost coloring of
the connected star.  Shared component vertices cannot be duplicated as
if they were independent channels.

\paragraph{Reversal and orientation metadata.}
Reversing a tile means interchanging its input and output lists and
reversing every component word.  It replaces $\sigma$ by $\sigma^{-1}$
and preserves the alignment condition.  The same theorem applies to
that explicitly reversed tile.  Edge directions or fibre-orientation
metadata do not affect an undirected cut cost, but any such data needed
for geometric provenance must still be transported separately.  A
negative fibre sign is not an instruction to complement a shore color.
If a different formalism imposes color-complementing seams or accesses
hidden corners, it is outside the vertex-identification model used here.

\paragraph{Sharpness and bounds not claimed.}
The cutoff four cannot be improved uniformly in this class, because the
class contains every connected tile.  For a one-edge tile, a path with
both endpoint colors zero needs at least four edges to have both marks:
the coloring $0,0,1,1,0$ realizes cost two on four edges, while three
edges cannot realize both monochromatic edge marks with those endpoints.
This establishes sharpness of the \emph{connected primitive cutoff},
not optimality of the factor $dQ$ in the routed size bound.

\subsection{Relation to existing graph-sparsification results}

The preservation of terminal cut information under graph replacement is
classical.  Hagerup, Katajainen, Nishimura, and Ragde developed mimicking
networks for external flow behavior \cite{hagerup1998}.  For thresholded
terminal connectivity, Chalermsook and coauthors construct
connectivity-$c$ mimicking networks of $O(pc^4)$ edges for $p$ terminals,
with algorithms polynomial in the explicit graph size for fixed $c$
\cite{chalermsook2021}.  Liu subsequently obtained
$\widetilde O(pc^3)$ edges with polynomial construction time in the
explicit graph size and $c$ \cite{liu2023}.

These results are essential context: this work does not claim that small-cut
preservation or polynomial-size terminal cut representatives are new.
The statement established here is a particular \emph{marked},
binary-multiplicity graph-grammar reduction, constructed without expanding
the input power, together with source-coloring lifting.  Ordinary terminal
minimum-cut values do not automatically determine the exact edge-mark
states used by the terminal pattern query.  The literature search located
no identical cyclewise formulation, but this is not a priority claim.

% Bibliography records for merging into the parent manuscript.
% The existing incoming report should be identified by its exact source hash
% in the final source audit; do not invent a repository commit for it.
